% ============================================================================
% categories/LATAM-PUB.tex — CATEGORY ARCHETYPE: Latin American public
%                            (federal/national) university
% ----------------------------------------------------------------------------
% LAYER 1 (category archetype) of the suite localization model. This file is
% \input AFTER shared/localization-defaults.tex (Layer 0) and BEFORE an
% institution-level set (Layer 2). It overrides the subset of localization keys
% that the "Latin American public university" archetype binds to a concrete
% generic value; every key it does NOT set keeps its Layer-0 default.
%
% ----------------------------------------------------------------------------
% WHICH UNIVERSITIES THIS ARCHETYPE FITS
% ----------------------------------------------------------------------------
% A research-active PUBLIC university in Latin America that is funded and
% supervised by the State (a federal or national ministry of education),
% holds the legal form of a public-law entity (in Brazil an autarquia or
% fundacao publica of special regime; the Spanish-speaking analogue is the
% autonomous national/state university), is governed by an elected rector
% with a collegial university council, and reaches the academic Internet
% through a State-backed national research-and-education network that is in
% turn federated regionally (RedCLARA) and globally (eduGAIN). It assumes a
% PUBLIC-SECTOR, BUDGET-CONSTRAINED estate: shared and NREN-provided
% infrastructure is the norm rather than the exception, and data custody is
% predominantly national/on-soil. Worked regulatory instance: BRAZIL (LGPD,
% E-Ciber, CTIR.Gov, the federal autarquia model, RNP as the NREN, CAFe as the
% identity federation). It generalises to other Latin American national public
% universities, whose RedCLARA-member NRENs and emerging data-protection
% regimes occupy the same structural roles.
%
% WHAT THIS FILE DELIBERATELY DOES NOT SET
%   * The ~3 country-specific NAME keys --- the specific national NREN, the
%     national data-protection authority, and the national/governmental CSIRT
%     --- are LEFT at their Layer-0 generic default. They are filled by a
%     Layer-2 institution-level set (e.g. RNP / ANPD / CTIR.Gov for a Brazilian
%     adopter, or the corresponding national bodies for another LATAM country),
%     because they vary by COUNTRY within this regional archetype.
%   * The regulatory keys below bind to the BRAZIL (BR) regime profile of the
%     shared Regulatory Applicability module
%     (shared/regulatory-applicability-EN.tex, "Brazil --- light regime"),
%     stated in ROLE terms so a non-Brazilian LATAM adopter can re-point them
%     without forking this file.
%
% Keys are kebab-case. Requires localization.sty (\setloc / \loc) and a prior
% \input of shared/localization-defaults.tex.
% ============================================================================

% ============================================================================
% STRUCTURAL OVERRIDES
% ============================================================================

% --- Subject institution & profile ------------------------------------------
\setloc{institution-type}{a State-funded public research university (a public-law entity supervised by the national ministry of education)}
\setloc{staffing-profile}{a public university operating under a State-appropriated budget, where IT capacity is constrained and shared/NREN-provided infrastructure is the default rather than the exception}

% --- Governance ladder ------------------------------------------------------
%     Public LATAM model: an elected rector heads the executive (reitoria /
%     rectorado); a collegial university council is the supreme deliberative
%     organ; the State ministry of education supervises the whole.
\setloc{executive-leadership}{the rectorate (the elected rector and the central executive administration of the public university)}
\setloc{governing-board}{the university council (the supreme collegial deliberative body) acting under the supervision of the national ministry of education}
\setloc{academic-governing-bodies}{the collegial academic councils (university, teaching-and-research, and faculty/unit councils)}

% --- Constituent units ------------------------------------------------------
\setloc{cs-research-unit}{a computing/informatics institute or department of the public university}
\setloc{security-research-centre}{a public-university research group or laboratory focused on cybersecurity and information security}
\setloc{institution-hpc-facility}{the university's shared high-performance-computing facility / team (often co-funded through the NREN or a regional consortium)}

% --- Pilot units ------------------------------------------------------------
\setloc{pilot-entities}{the university's shared HPC facility / team together with its computing/informatics institute}

% --- Operational roles ------------------------------------------------------
\setloc{internal-emergency-incident-contact}{the public university's information-security office / CSIRT as the internal incident-reporting contact, escalating to the national/governmental CSIRT}

% --- NREN tier & its services -----------------------------------------------
%     State-backed national NREN, regionally federated via RedCLARA, globally
%     via eduGAIN. Role-stated so a Brazilian adopter reads RNP and another
%     LATAM adopter reads its own RedCLARA-member NREN.
\setloc{nren-name}{the State-backed national research and education network (the country's RedCLARA-member NREN)}
\setloc{nren-csirt}{the national NREN's security team / CSIRT}
\setloc{peer-nren-roster}{peer Latin American NRENs federated through RedCLARA, each operating a national academic backbone, a CSIRT, and federated identity and shared-infrastructure services}
\setloc{nren-service-catalogue}{the national NREN's federated-service catalogue (federated identity, eduroam, video/collaboration, security monitoring and managed research-data transfer)}
\setloc{nren-research-cloud}{a national-NREN or RedCLARA-provided shared research cloud / e-infrastructure, where one exists}
\setloc{research-backbone-scale}{the regional research backbone (RedCLARA membership and backbone capacity across Latin America), with campus connections in the 1--100\,Gbps range subject to national funding}

% --- Identity federation ----------------------------------------------------
\setloc{identity-federation-profile}{the national academic identity federation (a RedCLARA-region federation interoperating with eduGAIN)}

% --- HPC tier ---------------------------------------------------------------
\setloc{national-hpc-facility}{a national or regional academic supercomputer (a State-funded national HPC centre or a RedCLARA/SCALAC-coordinated regional facility)}

% --- Deployment-scale model & data custody ----------------------------------
\setloc{deployment-custody-model}{the public university operates a shared or NREN-hosted deployment of its core components, or procures it under equivalent national-custody guarantees, sizing the estate to a constrained State-appropriated budget}
\setloc{data-custody-sovereignty-profile}{a national/on-soil custody arrangement (institution-, NREN- or domestic-provider-hosted) that keeps personal data under the national data-protection regime and minimises cross-border disclosure-regime exposure}

% ============================================================================
% REGULATORY OVERRIDES — BRAZIL (BR) regime profile
% ----------------------------------------------------------------------------
% Bound to the "Brazil --- light regime" row of the shared Regulatory
% Applicability module. Stated in role terms; a non-Brazilian LATAM adopter
% re-points these to its own jurisdiction via the module's 7-step method.
% NOTE: the country-specific NAME keys (national-dpa, national-gov-csirt,
% national NREN name) are intentionally NOT set here — see header.
% ============================================================================

\setloc{regulatory-regime}{binding national law in the adopter's Latin American jurisdiction (worked instance: the Brazilian federal stack — LGPD and the national cyber framework)}
\setloc{cybersecurity-baseline-regime}{the national cybersecurity framework carried, absent a single omnibus cyber statute, by the security limb of the national data-protection law and the national cyber strategy (worked instance: Brazil's LGPD Art.~46 with the E-Ciber strategy)}
\setloc{data-protection-regime}{the national general data-protection law (worked instance: Brazil's LGPD, Lei n\textsuperscript{o}~13.709/2018)}
\setloc{eid-trust-services-regime}{the national electronic-signature / trust-services and digital-inclusion regime (worked instance: Brazil's ICP-Brasil and the e-MAG accessibility standard under the Inclusion Law)}
\setloc{ai-regulation}{the adopter's national AI-governance regime as it emerges (worked instance: Brazil's pending AI bill, PL~2338/2023)}
\setloc{cross-border-transfer-regime}{the national international-transfer regime (self-assessed adequacy / safeguards / consent, with no EU-style adequacy list; worked instance: LGPD Chapter~V supervised by the national DPA)}
\setloc{essential-entity-classification-basis}{public universities as public-administration bodies in scope of the national data-protection and federal cyber-incident regimes (worked instance: federal universities as autarquias / fundacoes publicas of the federal public administration)}

\endinput
